\section{Resultados esperados en la configuración de la red IoT}

Los sensores estarán integrados en dispositivos diseñados específicamente para conectarse a una red IoT mediante protocolos de comunicación eficientes como LoRa o Wi-Fi. Estos dispositivos permitirán la transmisión de datos en tiempo real hacia una plataforma central, optimizando la cobertura y el consumo energético según las condiciones del entorno.

\begin{enumerate}[label=\Roman*.]
    \item Monitoreo en tiempo real de las concentraciones de gases de efecto invernadero, como \ch{CO2} y \ch{CH4}.
    \item Ampliación de la cobertura de monitoreo gracias al uso de LoRa en zonas rurales o remotas donde la conectividad Wi-Fi es limitada.
    \item Optimización del consumo energético de los dispositivos, aumentando su autonomía y reduciendo la necesidad de mantenimiento frecuente.
    \item Obtención de datos precisos y confiables mediante sensores calibrados y validados con métodos tradicionales.
    \item Implementación de análisis predictivo mediante inteligencia artificial para anticipar tendencias de emisiones.
    \item Escalabilidad del sistema mediante la incorporación de más sensores sin afectar el rendimiento de la red.
    \item Gestión eficiente y segura de los datos recolectados, con almacenamiento en la nube y visualización en dashboards.
    \item Contribución a la mitigación del cambio climático mediante la identificación oportuna de fuentes de emisión y la mejora de procesos agrícolas.
\end{enumerate}
